\section{Local convergence to the tree}\label{sec:localconv}

This section compares the neighbourhood of an agent in the near-tight hypergraph $H_q(K)$
with the corresponding ball of the Poisson hypertree $\tree$ of Section~\ref{sec:tree}. The
comparison is in total variation and the bound is explicit
(Theorem~\ref{thm:localconv}). It is the only fact about the finite market, apart from the
finiteness of $C_0$, that the tree bound of Section~\ref{sec:lp} uses. The leading error
term is $\tfrac9{10}K^2q^{-1/3}$ for each team of the ball, and it comes from near-tight
teams that share a pair.

\subsection{Balls, shapes and the two laws}\label{sec:localconv-setting}

Throughout the section $q\ge1$ and $r\ge1$ are integers and $K$ is a real number with
$0<K\le q^{2/3}$. We put
\[
s=Kq^{-2/3}\le1,\qquad \lambda=\frac{K^3}{6}=\Lam([0,K]),\qquad
S_r=\sum_{j=0}^{r-1}(2\lambda)^j .
\]
A team $e$ is \emph{near-tight} if $\cost_e\le K$, that is, if $\Defi(e)\le s$. The
condition $s\le1$ makes \eqref{eq:model-oneteam} available for every $t\le s$.

\begin{definition}\label{def:localconv-ball}
Let $G$ be a hypergraph whose vertices are agents, each with its sector, and whose edges are
teams, each marked by a cost. The \emph{distance} between two agents of $G$ is the least
number of teams in a chain joining them. For an agent $v$ the \emph{ball} $B_r(G,v)$ is the
rooted marked hypergraph formed by the root $v$, all teams of $G$ that contain an agent at
distance at most $r-1$ from $v$, the agents of these teams, their sectors and the costs. It
is taken up to bijections of the agents that preserve the root, the sectors, the teams and
the costs. The agents at distance at most $r-1$ from $v$ are the \emph{interior} agents of
the ball. For the near-tight hypergraph we write $B_r(o)=B_r(H_q(K),o)$.
\end{definition}

Let $\mathcal G_r$ be the set of finite rooted marked hypergraphs of this kind with costs in
$[0,K]$, up to the same bijections. An element of $\mathcal G_r$ is an unmarked rooted
hypergraph with sectors, of which there are countably many up to isomorphism, together with
a vector of costs modulo the finite automorphism group of the unmarked hypergraph. So
$\mathcal G_r$ is a countable disjoint union of quotients of cubes $[0,K]^n$ by finite
groups, and it carries the corresponding Borel $\sigma$-field. For probability measures
$\mu$, $\nu$ on a measurable space, $d_{\mathrm{TV}}(\mu,\nu)=\sup_{\mathcal A}
|\mu(\mathcal A)-\nu(\mathcal A)|$.

\emph{Tree shapes.} A \emph{tree shape} of radius $r$ with root sector $X$ is a finite
rooted hypertree $t$ of the following kind. The root has sector $X$. Every agent of depth at
most $r-1$ has a finite, linearly ordered set of child teams. Every team consists of its
parent agent and two further agents, one in each of the two other sectors. Agents of depth
$r$ have no child teams. There are countably many tree shapes. For a tree shape $t$ we write
$E(t)$ for its set of teams, $|E|=|E(t)|$ for their number, $V=1+2|E|$ for the number of
agents, and
\begin{align*}
m&=\#\{\text{agents of depth at most } r-1\},\\
P_2&=\#\{\text{unordered pairs of distinct teams with a common agent}\}.
\end{align*} Let $D_t\subseteq[0,K]^{E(t)}$ be the set of cost vectors
with pairwise distinct coordinates that increase along the child teams of every agent.

Let $\mathcal T_r\subseteq\mathcal G_r$ be the set of balls that are hypertrees (the
incidence graph of agents and teams is a tree) with pairwise distinct costs. In such a ball
the distance from the root is the depth, every team contains an interior agent, and so
agents of depth $r$ have no child teams. Ordering the child teams of every agent by
increasing cost, and telling the two further agents of a team apart by their sectors, we see
that every element of $\mathcal T_r$ with root sector $X$ has exactly one representation
$(t,\xi)$ with $t$ a tree shape and $\xi\in D_t$. We use two consequences. First, a
bijection of the agents of $(t,\xi)$ that preserves the root, the sectors, the teams and the
costs is the identity: it fixes the root, hence each child team of the root, since their
costs are distinct, hence the two further agents of each such team, since their sectors
differ, and so on by induction on the depth. Secondly, $\mathcal T_r$ carries the reference
measure $\varrho$ whose restriction to each shape $t$ is Lebesgue measure on $D_t$. For a
tree shape $t$ and a Borel set $A\subseteq D_t$ we write $[t,A]\subseteq\mathcal T_r$ for the
set of balls with representation $(t,\xi)$, $\xi\in A$. Every measurable subset of
$\mathcal T_r$ consisting of balls with root sector $X$ is a countable disjoint union of
such sets.

\emph{The two laws.} Let $X$ be a sector.
\begin{itemize}
\item $\mu=\mu_{q,K,r,X}$ is the law of $B_r(o)$ for an agent $o$ of sector $X$ of the
market. By symmetry it does not depend on the choice of $o$.
\item Let $\tree_K$ be the component of the root $o$ of the Poisson hypertree $\tree$
(Definition~\ref{def:tree-poisson}) after all teams of cost above $K$ have been deleted, as
in Remark~\ref{rem:tree-balls}, and let $\nu=\nu_{K,r,X}$ be the law of $B_r(\tree_K,o)$
when the root has sector $X$. By Definition~\ref{def:tree-poisson} and the restriction
theorem for Poisson processes, $\nu$ is the law of the following random marked
hypertree: the root has sector $X$, and every agent of depth at most $r-1$ has, independently
of everything else, child teams whose costs are the points of a Poisson process on $[0,K]$
with intensity $\xi^2/2\dd\xi$, of total mass $\lambda$; every child team has two new
agents, one in each of the two other sectors.
\end{itemize}

\subsection{The tree ball}\label{sec:localconv-treeball}

The facts about Poisson processes used in this section (the restriction theorem, the
density of the ordered points of a Poisson process of finite total mass, Campbell's formula)
are standard; see \cite{LastPenrose17}.

\begin{lemma}\label{lem:localconv-tree}
\begin{enumerate}
\item[\textup{(i)}] $\nu(\mathcal T_r)=1$, and $\nu$ has the density
\begin{equation}\label{eq:localconv-treedensity}
g(t,\xi)=e^{-m(t)\lambda}\prod_{e\in E(t)}\frac{\xi_e^2}{2}
\end{equation}
with respect to $\varrho$.
\item[\textup{(ii)}] For every measurable $\varphi\colon[0,K]\to[0,\infty)$,
\[
\E_\nu\sum_{e\in E}\varphi(\xi_e)=S_r\int_0^K\varphi(\xi)\,\frac{\xi^2}{2}\dd\xi .
\]
In particular $\E_\nu|E|=\lambda S_r$. Moreover $\E_\nu m=S_r$.
\item[\textup{(iii)}] $\E_\nu P_2=\dfrac{\lambda^2}{2}+\Bigl(\lambda+\dfrac{\lambda^2}{2}\Bigr)
(S_r-1)\le\Bigl(\lambda+\dfrac{\lambda^2}{2}\Bigr)S_r$.
\item[\textup{(iv)}] $Q_r:=\E_\nu V^2\le(r+1)^2(2r+1)\max(1,2\lambda)^{2r}$.
\end{enumerate}
\end{lemma}

\begin{proof}
(i) The points $\xi_1<\dots<\xi_n$ of a Poisson process on $[0,K]$ with intensity
$\xi^2/2\dd\xi$, listed in increasing order, have density $e^{-\lambda}\prod_i\xi_i^2/2$ on
$\{n\}\times\{\xi_1<\dots<\xi_n\}$; there is no factorial because the points are ordered.
The processes at the $m$ agents of depth at most $r-1$ are independent, which gives
\eqref{eq:localconv-treedensity} on the set where the child costs increase at every agent.
The costs are almost surely pairwise distinct, so $\nu(\mathcal T_r)=1$.

(ii) Let $Z_j$ be the number of agents of depth $j$, so $Z_0=1$. The variable $Z_j$ is a
function of the Poisson processes attached to agents of depth less than $j$, and is
independent of those attached to agents of depth $j$. Every agent has a Poisson($\lambda$)
number of child teams and every team has two new agents, so $\E Z_j=(2\lambda)^j$. By
Campbell's formula each agent of depth $j\le r-1$ contributes
$\int_0^K\varphi(\xi)\xi^2/2\dd\xi$ to the expectation of the sum over its child teams.
Summing over $j\le r-1$ gives the first formula; $\varphi\equiv1$ gives $\E_\nu|E|=\lambda
S_r$; and $m=\sum_{j<r}Z_j$ gives $\E_\nu m=S_r$.

(iii) Agents of depth $r$ lie in one team of the ball. The root, with $N$ child teams, lies
in $N$ teams and contributes $\binom N2$ pairs; any other agent of depth at most $r-1$ lies
in its parent team and $N$ child teams and contributes $\binom{N+1}2$. For a
Poisson($\lambda$) variable $N$, $\E\binom N2=\lambda^2/2$ and
$\E\binom{N+1}2=\lambda+\lambda^2/2$. The number of child teams of an agent is independent
of the presence of the agent, and the expected number of agents of depth between $1$ and
$r-1$ is $S_r-1$.

(iv) $(Z_j)$ is a Galton--Watson process whose offspring number is twice a
Poisson($\lambda$) variable, with mean $2\lambda$ and variance $4\lambda$. The conditional
variance formula gives $\Var Z_j=(2\lambda)^2\Var Z_{j-1}+4\lambda(2\lambda)^{j-1}$, hence
\[
\E Z_j^2=(2\lambda)^{2j}+2\sum_{i=0}^{j-1}(2\lambda)^{j+i}\le(1+2j)\max(1,2\lambda)^{2j}.
\]
Since $V=\sum_{j=0}^rZ_j$, the Cauchy--Schwarz inequality gives
$\E V^2\le(r+1)\sum_{j=0}^r\E Z_j^2\le(r+1)^2(2r+1)\max(1,2\lambda)^{2r}$.
\end{proof}

\subsection{One team and its neighbours}\label{sec:localconv-oneteam}

\begin{lemma}\label{lem:localconv-oneteam}
Let $e$ be a team with pair deficits $\defi_1,\defi_2,\defi_3$.
\begin{enumerate}
\item[\textup{(i)}] The cost $\cost_e$ has density $\xi^2/(2q^2)$ on $[0,K]$, and
$\P(e\text{ is near-tight})=s^3/6=\lambda/q^2$. Conditionally on $e$ being near-tight,
$\cost_e$ has density $3\xi^2/K^3$ on $[0,K]$, the vector $(\defi_1,\defi_2,\defi_3)$ is
uniform on the solid simplex $\{x\in[0,\infty)^3:\ x_1+x_2+x_3\le s\}$, and $\defi_1$ has
density $3(s-x)^2/s^3$ on $[0,s]$.
\item[\textup{(ii)}] Let $\delta',\delta''$ be independent and uniform on $[0,1]$. For
$x\in[0,s]$, $\P(x+\delta'+\delta''\le s)=(s-x)^2/2$. For $x,x'\ge0$,
$\P(x+x'+\delta'\le s)\le s$.
\item[\textup{(iii)}] $\E\bigl[(s-\defi_1)^2/2\bigm|e\text{ is near-tight}\bigr]=3s^2/10$.
Consequently, if $\psi(\xi)=\E[(s-\defi_1)^2/2\mid\cost_e=\xi]$ for $\xi\in[0,K]$, then
\begin{equation}\label{eq:localconv-psi}
\int_0^K\psi(\xi)\,\frac{3\xi^2}{K^3}\dd\xi=\frac{3s^2}{10}.
\end{equation}
By symmetry $\psi$ is the same for $\defi_2$ and $\defi_3$.
\item[\textup{(iv)}] \emph{(Pair-sharing rate.)} Given the pair deficits of $e$, the
$3(q-1)$ teams that share a pair with $e$ are near-tight independently of each other, the
one through the pair $p$ of $e$ with probability $(s-\defi(p))^2/2$ when $e$ is near-tight.
Hence,
if $\pi(e)$ denotes the number of near-tight teams other than $e$ that share a pair with
$e$,
\begin{equation}\label{eq:localconv-pairsharing}
\E\bigl[\pi(e)\bigm|e\text{ is near-tight}\bigr]
=3(q-1)\cdot\frac{3s^2}{10}=\frac9{10}\,K^2q^{-1/3}\Bigl(1-\frac1q\Bigr).
\end{equation}
\end{enumerate}
\end{lemma}

\begin{proof}
(i) Since $s\le1$, for $t\le s$ the simplex $\{x\ge0,\ x_1+x_2+x_3\le t\}$ lies in the unit
cube and has volume $t^3/6$; this is \eqref{eq:model-oneteam}. The density of
$\cost_e=q^{2/3}\Defi(e)$ follows by differentiation, and the conditional statements because
the deficits are uniform on the cube. The section of the simplex at $x_1=x$ has area
$(s-x)^2/2$, and $\bigl((s-x)^2/2\bigr)/(s^3/6)=3(s-x)^2/s^3$.

(ii) The triangle $\{\delta',\delta''\ge0,\ \delta'+\delta''\le s-x\}$ lies in the unit
square and has area $(s-x)^2/2$. The second event is contained in $\{\delta'\le s\}$.

(iii) By (i), $\int_0^s\frac{(s-x)^2}2\cdot\frac{3(s-x)^2}{s^3}\dd x=\frac{3}{2s^3}\cdot
\frac{s^5}{5}=\frac{3s^2}{10}$. Formula \eqref{eq:localconv-psi} is the tower property,
because $\cost_e$ has density $3\xi^2/K^3$ given that $e$ is near-tight.

(iv) A team $\alpha\ne e$ that shares a pair with $e$ shares exactly one pair $p$ with it,
and there are $q-1$ such teams for each of the three pairs of $e$. The other two pairs of
$\alpha$ are not pairs of $e$, and for different $\alpha$ these pairs are different. So,
given the deficits of $e$, the events $\{\alpha\text{ is near-tight}\}$ are independent,
with the probabilities given by (ii). Take expectations with (iii).
\end{proof}

\subsection{The theorem}\label{sec:localconv-theorem}

\begin{theorem}\label{thm:localconv}
Let $q\ge1$, $r\ge1$, $0<K\le q^{2/3}$, and let $X$ be a sector. Then
\begin{equation}\label{eq:localconv-bound}
\begin{aligned}
d_{\mathrm{TV}}\bigl(\mu_{q,K,r,X},\nu_{K,r,X}\bigr)&\le\eps_1(q,K,r),\\
\eps_1(q,K,r)&:=\frac9{10}\,K^2q^{-1/3}\cdot\lambda S_r
+(2\lambda+\lambda^2)\,S_r\,Kq^{-2/3}\\
&\qquad+\frac{Q_r}{q}
+\frac{\lambda^2S_r}{q^2},
\end{aligned}
\end{equation}
where $\lambda=K^3/6$, $S_r=\sum_{j<r}(2\lambda)^j$ and
$Q_r\le(r+1)^2(2r+1)\max(1,2\lambda)^{2r}$ is as in Lemma~\ref{lem:localconv-tree}. The four
terms account, in this order, for a near-tight team that shares a pair with one team of the
ball; a near-tight team that shares pairs with two teams of the ball; coincidences among the
agents when the ball is placed in the market; and the replacement of a binomial count by a
Poisson one. In particular, for an agent $o$ of the market,
\[
\P\bigl(B_r(o)\text{ is not a hypertree}\bigr)\le\eps_1(q,K,r),
\]
and $\eps_1(q,K,r)\to0$ as $q\to\infty$ for fixed $K$ and $r$. The same bound holds for the
total variation distance between the law of $B_r(H_q(K),v)$, where $v$ is a uniformly random
agent independent of the market, and the law of $B_r(\tree_K,o)$, where the root sector of
$\tree$ is uniform.
\end{theorem}

The first term is the pair-sharing rate \eqref{eq:localconv-pairsharing} for one team,
multiplied by the expected number $\lambda S_r$ of teams in the ball of the tree.

\begin{proof}
Fix an agent $o$ of sector $X$, a tree shape $t$ with root sector $X$ and a Borel set
$A\subseteq D_t$.

\emph{Step 1: the event $\{B_r(o)\in[t,A]\}$.} An \emph{embedding} is an injective map
$\phi$ from the agents of $t$ to the agents of the market that preserves sectors and sends
the root to $o$. For a team $e$ of $t$, $\phi(e)$ is the team formed by the images of its
agents; these are the \emph{tree teams} of $\phi$, and the pairs contained in them are its
\emph{tree pairs}. Put $\xi_\phi=(\cost_{\phi(e)})_{e\in E(t)}$, and let $W_\phi$ be the number
of near-tight teams that are not tree teams and contain $\phi(a)$ for some agent $a$ of $t$
of depth at most $r-1$. We claim that
\begin{equation}\label{eq:localconv-union}
\{B_r(o)\in[t,A]\}=\bigsqcup_{\phi}\ \{\xi_\phi\in A\}\cap\{W_\phi=0\},
\end{equation}
a disjoint union over all embeddings. Suppose $\xi_\phi\in A$ and $W_\phi=0$. Then the tree
teams are near-tight, and by induction on $j\le r-1$ the agents at distance $j$ from $o$ in
$H_q(K)$ are exactly the images of the agents of depth $j$, because every near-tight team
that contains such an image is a tree team. Hence the teams of $B_r(o)$ are exactly the tree
teams, and $B_r(o)$ has the representation $(t,\xi_\phi)$. Conversely, if $B_r(o)\in[t,A]$,
an isomorphism from $(t,\xi)$ onto $B_r(o)$ is an embedding $\phi$ with $\xi_\phi=\xi\in A$,
and every near-tight team containing an interior agent of $B_r(o)$ belongs to $B_r(o)$, so
$W_\phi=0$. The union is disjoint because two embeddings for which the event holds differ by
a bijection of the agents of $(t,\xi)$ that preserves the root, the sectors, the teams and
the costs, and such a bijection is the identity.

Two different tree teams share at most one agent, because $t$ is a hypertree. So the $3|E|$
tree pairs are distinct, the triples of pair deficits of the tree teams are independent, and
by Lemma~\ref{lem:localconv-oneteam}(i) the vector $\xi_\phi$ has density
$\prod_e\xi_e^2/(2q^2)$ on $[0,K]^{E(t)}$. Write $\mathbf d_\phi$ for the family of deficits
of the tree pairs. Then, by \eqref{eq:localconv-union},
\begin{equation}\label{eq:localconv-mass}
\mu([t,A])=\sum_\phi\E\bigl[\ind\{\xi_\phi\in A\}\,\P(W_\phi=0\mid\mathbf d_\phi)\bigr].
\end{equation}
The agents of $t$ other than the root are placed one at a time, and at each of the $2|E|$
steps at least $q-V$ agents of the required sector are free. So the number of embeddings is
at least
\begin{equation}\label{eq:localconv-count}
(q-V)_+^{2|E|}\ \ge\ q^{2|E|}\Bigl(1-\frac{2|E|V}{q}\Bigr)_+ ,
\end{equation}
by Bernoulli's inequality when $V\le q$; when $V>q$ the right-hand side is $0$.

\emph{Step 2: the other teams.} Fix an embedding $\phi$ and let $F$ be the set of teams that
are not tree teams and contain the image of an agent of depth at most $r-1$, so that
$W_\phi$ is the number of near-tight teams in $F$. Classify the teams of $F$ by the number of
tree pairs they contain.
\begin{itemize}
\item $F_0$: no tree pair. Each agent lies in $q^2$ teams, so $|F_0|\le mq^2$. Given
$\mathbf d_\phi$, a team of $F_0$ is near-tight with probability $s^3/6=\lambda/q^2$.
\item $F_1$: exactly one tree pair $p$, contained in the tree team $\phi(e)$. The team is one
of the $q-1$ other teams through $p$, and its other two pairs are not tree pairs. Given
$\mathbf d_\phi$, on the event $\{\xi_\phi\in A\}$ (where $\defi(p)\le s$) it is near-tight
with probability $(s-\defi(p))^2/2$, by Lemma~\ref{lem:localconv-oneteam}(ii).
\item $F_2$: exactly two tree pairs. Two pairs of a team have a common agent $b$. If both
were contained in the same tree team, the team would be that tree team. So they lie in two
different tree teams $\phi(e_1)$, $\phi(e_2)$ through $b$, and the team consists of $b$, one
further agent of $\phi(e_1)$ and the agent of $\phi(e_2)$ in the remaining sector. There are
two such teams for each pair $\{e_1,e_2\}$ of teams of $t$ with a common agent, so
$|F_2|\le2P_2$. The third pair of the team is not a tree pair, so given $\mathbf d_\phi$ the
team is near-tight with probability at most $s$, by Lemma~\ref{lem:localconv-oneteam}(ii).
\item Three tree pairs: impossible. No two of them lie in the same tree team, so they would
lie in three different tree teams that pairwise share an agent of the team, and this is a
cycle of length six in the incidence graph of $\phi(t)$.
\end{itemize}

\emph{Step 3: Harris' inequality.} Conditionally on $\mathbf d_\phi$, the deficits of all
other pairs are independent and uniform on $[0,1]$, and for each $\alpha\in F$ the event
$\{\alpha\text{ is not near-tight}\}=\{\Defi(\alpha)>s\}$ is increasing in them. By Harris'
inequality \cite{Harris60}, with $p_\alpha=\P(\alpha\text{ is near-tight}\mid\mathbf d_\phi)$,
\[
\P(W_\phi=0\mid\mathbf d_\phi)\ \ge\ \prod_{\alpha\in F}(1-p_\alpha).
\]
On the event $\{\xi_\phi\in A\}$ put
\[
Y_\phi=(q-1)\sum_{e\in E(t)}\ \sum_{p\subset\phi(e)}\frac{(s-\defi(p))^2}{2},
\]
the inner sum being over the three pairs of $\phi(e)$. By Step~2,
$\sum_{\alpha\in F_1}p_\alpha\le Y_\phi$ and $\sum_{\alpha\in F_2}p_\alpha\le2P_2s$. We use
$\prod_i(1-p_i)\ge1-\sum_ip_i$ for $p_i\in[0,1]$ on $F_1\cup F_2$, and on $F_0$ the
inequality $\log(1-u)\ge-u-u^2$ with $u=\lambda/q^2\le1/6$ (recall $K\le q^{2/3}$), which
gives $(1-\lambda/q^2)^{mq^2}\ge e^{-m\lambda}(1-m\lambda^2/q^2)$. Hence, on
$\{\xi_\phi\in A\}$,
\begin{equation}\label{eq:localconv-harris}
\P(W_\phi=0\mid\mathbf d_\phi)\ \ge\ e^{-m\lambda}\Bigl(1-\frac{m\lambda^2}{q^2}\Bigr)_+
\bigl(1-Y_\phi-2P_2s\bigr)_+ .
\end{equation}
(The positive parts may be taken because the left-hand side is nonnegative.)

\emph{Step 4: a lower bound for the density of $\mu$ on $\mathcal T_r$.} The triples of pair
deficits of the tree teams are independent, so for a pair $p$ of $\phi(e)$ we have
$\E[(s-\defi(p))^2/2\mid\xi_\phi]=\psi(\cost_{\phi(e)})$, with $\psi$ as in
Lemma~\ref{lem:localconv-oneteam}(iii). Therefore, with
\[
\omega_t(\xi)=1-3(q-1)\sum_{e\in E(t)}\psi(\xi_e)-2P_2Kq^{-2/3},
\]
and since $s=Kq^{-2/3}$,
\begin{align*}
\E\bigl[\ind\{\xi_\phi\in A\}\,(1-Y_\phi-2P_2s)_+\bigr]
&\ge\E\bigl[\ind\{\xi_\phi\in A\}\,(1-Y_\phi-2P_2s)\bigr]\\
&=\int_A\omega_t(\xi)\prod_{e\in E(t)}\frac{\xi_e^2}{2q^2}\dd\xi .
\end{align*}
Insert \eqref{eq:localconv-harris} into \eqref{eq:localconv-mass}. Every term of the sum
over $\phi$ is then bounded below by the same nonnegative quantity, so the sum is at least
that quantity times the lower bound \eqref{eq:localconv-count} for the number of embeddings;
then apply the last display, whose left-hand side is multiplied by a nonnegative factor, and
compare with \eqref{eq:localconv-treedensity}. With
\begin{equation}\label{eq:localconv-ratio}
\underline f(t,\xi)=\Bigl(1-\frac{2|E|V}{q}\Bigr)_+\Bigl(1-\frac{m\lambda^2}{q^2}\Bigr)_+
\,\omega_t(\xi)
\end{equation}
we obtain $\mu([t,A])\ge\int_{[t,A]}\underline f\dd\nu$. Since the sets $[t,A]$ exhaust the
measurable subsets of $\mathcal T_r$ with root sector $X$ up to countable disjoint unions,
\begin{equation}\label{eq:localconv-lower}
\mu(\mathcal A)\ \ge\ \int_{\mathcal A}\underline f_+\dd\nu\qquad\text{for every measurable }
\mathcal A\subseteq\mathcal T_r
\end{equation}
(apply the bound to $\mathcal A\cap\{\underline f>0\}$). The first two factors of
\eqref{eq:localconv-ratio} lie in $[0,1]$ and the third is at most $1$, so
\begin{equation}\label{eq:localconv-oneminus}
(1-\underline f)_+\ \le\ \frac{2|E|V}{q}+\frac{m\lambda^2}{q^2}+2P_2Kq^{-2/3}
+3(q-1)\sum_{e\in E(t)}\psi(\xi_e).
\end{equation}
Indeed, if the third factor is negative then $\underline f$ is at least that factor; and if
it is nonnegative, $(1-a)_+(1-b)_+(1-c)\ge1-a-b-c$ for $a,b\ge0$ and $c\in[0,1]$.

\emph{Step 5: total variation.} Only the lower bound \eqref{eq:localconv-lower} is needed.
For measurable $\mathcal A\subseteq\mathcal G_r$, since $\nu(\mathcal T_r)=1$,
\[
\nu(\mathcal A)-\mu(\mathcal A)\le\nu(\mathcal A\cap\mathcal T_r)-\mu(\mathcal A\cap
\mathcal T_r)\le\int_{\mathcal A\cap\mathcal T_r}(1-\underline f_+)\dd\nu
\le\int(1-\underline f)_+\dd\nu ,
\]
and $d_{\mathrm{TV}}(\mu,\nu)=\sup_{\mathcal A}(\nu(\mathcal A)-\mu(\mathcal A))$ for
probability measures. It remains to integrate \eqref{eq:localconv-oneminus} against $\nu$,
with Lemma~\ref{lem:localconv-tree}.
\begin{itemize}
\item By Lemma~\ref{lem:localconv-tree}(ii) with $\varphi=\psi$ and by
\eqref{eq:localconv-psi},
\[
\E_\nu\sum_e\psi(\xi_e)=S_r\int_0^K\psi(\xi)\,\frac{\xi^2}2\dd\xi=\lambda S_r\cdot
\frac{3s^2}{10}.
\]
So the last term has $\nu$-mean
$3(q-1)\cdot\frac{3s^2}{10}\cdot\lambda S_r\le\frac9{10}K^2q^{-1/3}\cdot\lambda S_r$.
\item By Lemma~\ref{lem:localconv-tree}(iii),
$\E_\nu\bigl[2P_2Kq^{-2/3}\bigr]\le(2\lambda+\lambda^2)S_rKq^{-2/3}$.
\item $2|E|V\le V^2$, so $\E_\nu[2|E|V/q]\le Q_r/q$.
\item $\E_\nu m=S_r$, so $\E_\nu[m\lambda^2/q^2]=\lambda^2S_r/q^2$.
\end{itemize}
The sum is $\eps_1(q,K,r)$. Every ball in $\mathcal T_r$ is a hypertree and
$\nu(\mathcal T_r)=1$, so $\P(B_r(o)\text{ is not a hypertree})\le\mu(\mathcal G_r\setminus
\mathcal T_r)\le d_{\mathrm{TV}}(\mu,\nu)\le\eps_1$. For fixed $K$ and $r$ every term of
$\eps_1$ tends to $0$. Finally, a uniformly random agent has a uniformly random sector, and
the total variation distance between two mixtures with the same weights is at most the
largest distance between the components.
\end{proof}

\subsection{Two roots}\label{sec:localconv-two}

\begin{proposition}\label{prop:localconv-two}
Let $o\ne o'$ be agents with sectors $X$, $X'$, and let $\mu_2$ be the joint law of
$(B_r(o),B_r(o'))$. Under the hypotheses of Theorem~\ref{thm:localconv},
\[
d_{\mathrm{TV}}\bigl(\mu_2,\ \nu_{K,r,X}\otimes\nu_{K,r,X'}\bigr)\ \le\
2\,\eps_1(q,K,r)+\frac{2Q_r}{q}\ \le\ 4\,\eps_1(q,K,r).
\]
\end{proposition}

\begin{proof}
Repeat the proof of Theorem~\ref{thm:localconv} for a pair of tree shapes $t,t'$ with root
sectors $X,X'$, a Borel set $A\subseteq D_t\times D_{t'}$ and pairs of embeddings
$(\phi,\phi')$ of $t$ and $t'$ with $\phi(\text{root})=o$, $\phi'(\text{root})=o'$ and
disjoint images. Quantities of $t'$ carry a prime. The tree teams are now those of $\phi$
and of $\phi'$, and $W_{\phi,\phi'}$ counts the near-tight teams that are not tree teams and
contain the image of an agent of depth at most $r-1$ of $t$ or of $t'$.

\emph{Step 1.} If $(\xi_\phi,\xi_{\phi'})\in A$ and $W_{\phi,\phi'}=0$, then the induction
of Step~1 above, started at $o$, stays inside the image of $\phi$, because a tree team that
contains an agent of the image of $\phi$ is a tree team of $\phi$. So $B_r(o)$ has the
representation $(t,\xi_\phi)$, and likewise $B_r(o')$ has the representation
$(t',\xi_{\phi'})$. The events for different pairs $(\phi,\phi')$ are disjoint as before.
Hence $\{(B_r(o),B_r(o'))\in[t,t',A]\}$ \emph{contains} the disjoint union of the events
$\{(\xi_\phi,\xi_{\phi'})\in A\}\cap\{W_{\phi,\phi'}=0\}$. It may be larger, since the two
balls may overlap; this is harmless, because only a lower bound for $\mu_2$ is needed. The
number of pairs of embeddings with disjoint images is at least
$q^{2|E|+2|E'|}\bigl(1-2(|E|+|E'|)(V+V')/q\bigr)_+$. Tree teams of the same embedding share
at most one agent and tree teams of different embeddings share none, so the $3(|E|+|E'|)$
tree pairs are distinct and the costs of the tree teams are independent.

\emph{Step 2.} $|F_0|\le(m+m')q^2$. $F_1$ is as before for each tree team. A team with
exactly two tree pairs contains their common agent $b$, and the two tree teams through $b$
belong to the same embedding; so $|F_2|\le2(P_2+P_2')$. Three tree pairs are impossible as
before.

\emph{Steps 3 to 5.} Harris' inequality applies as before, and the density of $\mu_2$ with
respect to $\nu_{K,r,X}\otimes\nu_{K,r,X'}$ on $\mathcal T_r\times\mathcal T_r$ is bounded
below by a function $\underline f_2$ with
\begin{align*}
(1-\underline f_2)_+\le{}&\frac{2(|E|+|E'|)(V+V')}{q}+\frac{(m+m')\lambda^2}{q^2}
+2(P_2+P_2')Kq^{-2/3}\\
&+3(q-1)\Bigl(\sum_{e\in E(t)}\psi(\xi_e)+\sum_{e\in E(t')}\psi(\xi'_e)\Bigr).
\end{align*}
The last three terms have mean at most twice the corresponding terms of $\eps_1$. For the
first, $2(|E|+|E'|)(V+V')\le(V+V')^2\le2V^2+2V'^2$, with mean at most $4Q_r$.
\end{proof}

\begin{corollary}\label{cor:localconv-prob}
Let $h\colon\mathcal G_r\to[-1,1]$ be measurable, let
\[
A_q=\frac1{3q}\sum_vh\bigl(B_r(H_q(K),v)\bigr)
\]
be its average over the $3q$ agents, and let $\bar\nu(h)$ be the mean of
$h(B_r(\tree_K,o))$ when the root sector is uniform. Then
\[
|\E A_q-\bar\nu(h)|\le2\,\eps_1(q,K,r),\qquad\Var A_q\le\frac1{3q}+12\,\eps_1(q,K,r).
\]
In particular $A_q\to\bar\nu(h)$ in probability as $q\to\infty$ for fixed $K$ and $r$.
\end{corollary}

\begin{proof}
If $|h|\le1$ and $\mu,\nu$ are probability measures, then
$|\int h\dd\mu-\int h\dd\nu|\le2\,d_{\mathrm{TV}}(\mu,\nu)$. This and
Theorem~\ref{thm:localconv} give the first bound. Let $o\ne o'$ be agents with sectors
$X,X'$, and write $\nu_X(h)=\int h\dd\nu_{K,r,X}$, $h_o=h(B_r(o))$. By
Proposition~\ref{prop:localconv-two} and Theorem~\ref{thm:localconv},
\[
\bigl|\E[h_oh_{o'}]-\nu_X(h)\nu_{X'}(h)\bigr|\le8\eps_1,\qquad
\bigl|\E h_o\,\E h_{o'}-\nu_X(h)\nu_{X'}(h)\bigr|\le4\eps_1 .
\]
So the covariance of $h_o$ and $h_{o'}$ is at most $12\eps_1$ in absolute value. The $3q$
diagonal terms of the variance of $\sum_vh_v$ are at most $1$ each.
\end{proof}

The corollary is not used in the sequel: Section~\ref{sec:lp} needs only expectations, that
is, Theorem~\ref{thm:localconv}. Note the rates for fixed $K$ and $r$: the bias of $A_q$ is
$O(q^{-1/3})$, and its fluctuation is $O(q^{-1/6})$ in probability.

\subsection{Remarks}\label{sec:localconv-remarks}

\begin{remark}[Reading the bound]\label{rem:localconv-reading}
For fixed $K$ and $r$ the first term of $\eps_1$ is of order $q^{-1/3}$ and the other three
are $O(q^{-2/3})$. For $K\ge2$ we have $2\lambda\ge8/3$ and $\lambda S_r\le\frac45
(2\lambda)^r$, so the first term is at most
\[
0.72\,K^2\Bigl(\frac{K^3}{3}\Bigr)^{r}q^{-1/3}.
\]
The bound is therefore informative at radius $r$ only when $q^{1/3}$ is large compared with
$K^{3r+2}/3^r$, and for a single team the requirement is that $q^{1/3}$ be large compared
with $K^2$. The theorem is for fixed $K$ and $r$. Nothing here controls $K$ growing with
$q$; by \eqref{eq:localconv-pairsharing}, when $K^2$ is comparable to $q^{1/3}$ a near-tight
team shares a pair with a number of near-tight teams that is of order one in expectation.
\end{remark}

\begin{remark}[Where pairs enter]\label{rem:localconv-pairs}
The fact that a near-tight team has small pair deficits, which favours the other teams
through its pairs, enters the proof only through the classes $F_1$ and $F_2$. It produces
the rate $K^2q^{-1/3}$. Teams that share only an agent with the ball contribute nothing
beyond the last term of $\eps_1$.
\end{remark}

\begin{remark}[Floating point: how large $q$ must be]\label{rem:localconv-size}
The numbers in this remark are floating-point evidence and are not used in any proof. The
identity \eqref{eq:localconv-pairsharing} is exact, and it shows that the tree regime
starts late: the expected number of near-tight partners of a near-tight team through a
shared pair is below $1/10$ only when $q$ is at least about $(9K^2)^3$, that is, $3.0\cdot10^6$ for
$K=4$ and $3.4\cdot10^7$ for $K=6$. Simulations of fresh markets with all $q^3$ teams
enumerated (script \texttt{lemmaC/round3/localconv/pairshare.py}, output
\texttt{out\_pairshare.txt} in the same folder; $300$, $80$ and $14$ markets at
$q=32,64,128$) agree with \eqref{eq:localconv-pairsharing} to within sampling error: at
$K=4$ the mean number of partners per near-tight team was $4.41$, $3.60$, $2.81$ against
$4.39$, $3.54$, $2.84$ from the formula. In the same runs, at $K=4$ the fraction of
near-tight teams with at least one such partner was $0.987$, $0.970$, $0.934$, and the
fraction of agents whose ball of radius $1$ is not a hypertree was $0.952$, $0.952$,
$0.945$; at $K=2$ the latter fraction was $0.235$, $0.210$, $0.185$. So at every size that
can be simulated the near-tight hypergraph is far from a hypertree, and
Theorem~\ref{thm:localconv}, like every bound that rests on it, says nothing at such sizes.
\end{remark}
